% GENERATED FILE. Edit canonical lessons, then run npm run generate:books.
% canonical-source-hash: fnv1a64:01874671d6269292
% canonical-lessons: BN-C175-durgondho, BN-C175-jharon, BN-C175-boti, BN-C175-shilnora, BN-C175-gs

\chapter{Stench, Duster, Bonti, Grinding stone, Gas}
\label{ch:bn-stench-duster-bonti-grinding-stone-gas}
\hlchaptermodality{175}

\begin{chapteropening}
\textbf{By the end of this chapter:} \emph{I can say a stench, a duster, a bonti, a grinding stone and gas.}

Complete the last lesson of chapter 175: I can say a stench, a duster, a bonti, a grinding stone and gas.
\end{chapteropening}

\section[\texorpdfstring{\bn{দুর্গন্ধ} (durgôndhô)}{দুর্গন্ধ (durgôndhô)}]{\bn{দুর্গন্ধ} (durgôndhô) — a stench}

\label{lesson:BN-C175-durgondho}



\begin{quote}
Before the new one: say the Bengali for a tyre, then the Bengali for a brake.
\end{quote}

\subsection*{You'll want to know: \bn{দুর্গন্ধ}}
\textbf{\bn{দুর্গন্ধ}} — \emph{durgôndhô} — \textquotedblleft{}a stench\textquotedblright{}.

\begin{cousinweb}
\textbf{\bn{দুর্গন্ধ}} (\emph{durgôndhô}) is Sanskrit \emph{durgandha}, \textquotedblleft{}bad smell\textquotedblright{}, taken over whole as a learned loan, a tatsama (CDIAL 6431).
\end{cousinweb}

\begin{grammarlens}[title={What you've built}]
One more word for this chapter.
\end{grammarlens}

\subsection*{Your turn}
\begin{itemize}
\raggedright
  \item \emph{Say it:} \emph{durgôndhô}
  \item \emph{Say it:} \emph{durgôndhô}, once more
  \item \emph{Say it:} say \emph{durgôndhô}
  \item \emph{Recall:} say the Bengali for a tyre, then the Bengali for a brake, then say \emph{durgôndhô} again
\end{itemize}

\begin{tcolorbox}[breakable,colback=teal!4,colframe=teal!35!black,title={Before you move on}]
What does \bn{দুর্গন্ধ} mean? (\textquotedblleft{}A stench\textquotedblright{}.) Say it once more.
\end{tcolorbox}
\section[\texorpdfstring{\bn{ঝাড়ন} (jhāṛôn)}{ঝাড়ন (jhāṛôn)}]{\bn{ঝাড়ন} (jhāṛôn) — a duster}

\label{lesson:BN-C175-jharon}



\begin{quote}
Before the new one: say the Bengali for a brake, then the Bengali for a stench.
\end{quote}

\subsection*{You'll want to know: \bn{ঝাড়ন}}
\textbf{\bn{ঝাড়ন}} — \emph{jhāṛôn} — \textquotedblleft{}a duster\textquotedblright{}.

\begin{cousinweb}
\textbf{\bn{ঝাড়ন}} (\emph{jhāṛôn}) is built on the verb \emph{jhāṛā}, \textquotedblleft{}to shake off\textquotedblright{}, which CDIAL lists under \emph{*jhāṭayati}, \textquotedblleft{}shakes down\textquotedblright{}, a form it reconstructs (CDIAL 5328), + \emph{-ôn}, an ending for a tool: \textquotedblleft{}a thing for shaking off dust\textquotedblright{}.
\end{cousinweb}

\begin{grammarlens}[title={What you've built}]
One more word for this chapter.
\end{grammarlens}

\subsection*{Your turn}
\begin{itemize}
\raggedright
  \item \emph{Say it:} \emph{jhāṛôn}
  \item \emph{Say it:} \emph{jhāṛôn}, once more
  \item \emph{Say it:} say \emph{jhāṛôn}
  \item \emph{Recall:} say the Bengali for a brake, then the Bengali for a stench, then say \emph{jhāṛôn} again
\end{itemize}

\begin{tcolorbox}[breakable,colback=teal!4,colframe=teal!35!black,title={Before you move on}]
What does \bn{ঝাড়ন} mean? (\textquotedblleft{}A duster\textquotedblright{}.) Say it once more.
\end{tcolorbox}
\section[\texorpdfstring{\bn{বঁটি} (bõṭi)}{বঁটি (bõṭi)}]{\bn{বঁটি} (bõṭi) — a bonti (a kitchen blade)}

\label{lesson:BN-C175-boti}



\begin{quote}
Before the new one: say the Bengali for a stench, then the Bengali for a duster.
\end{quote}

\subsection*{You'll want to know: \bn{বঁটি}}
\textbf{\bn{বঁটি}} — \emph{bõṭi} — \textquotedblleft{}a bonti (a kitchen blade)\textquotedblright{}.

\begin{cousinweb}
The origin of \textbf{\bn{বঁটি}} (\emph{bõṭi}) is not settled: CDIAL has no entry for it, and a search by form and by the glosses \textquotedblleft{}blade\textquotedblright{} and \textquotedblleft{}knife\textquotedblright{} finds only words of other shapes, such as Bengali \emph{phal}, \textquotedblleft{}blade\textquotedblright{} (CDIAL 9052).
\end{cousinweb}

\begin{grammarlens}[title={What you've built}]
One more word for this chapter.
\end{grammarlens}

\subsection*{Your turn}
\begin{itemize}
\raggedright
  \item \emph{Say it:} \emph{bõṭi}
  \item \emph{Say it:} \emph{bõṭi}, once more
  \item \emph{Say it:} say \emph{bõṭi}
  \item \emph{Recall:} say the Bengali for a stench, then the Bengali for a duster, then say \emph{bõṭi} again
\end{itemize}

\begin{tcolorbox}[breakable,colback=teal!4,colframe=teal!35!black,title={Before you move on}]
What does \bn{বঁটি} mean? (\textquotedblleft{}A bonti (a kitchen blade)\textquotedblright{}.) Say it once more.
\end{tcolorbox}
\section[\texorpdfstring{\bn{শিলনোড়া} (shilnoṛā)}{শিলনোড়া (shilnoṛā)}]{\bn{শিলনোড়া} (shilnoṛā) — a grinding stone}

\label{lesson:BN-C175-shilnora}



\begin{quote}
Before the new one: say the Bengali for a duster, then the Bengali for a bonti.
\end{quote}

\subsection*{You'll want to know: \bn{শিলনোড়া}}
\textbf{\bn{শিলনোড়া}} — \emph{shilnoṛā} — \textquotedblleft{}a grinding stone\textquotedblright{}.

\begin{cousinweb}
\textbf{\bn{শিলনোড়া}} (\emph{shilnoṛā}) is \emph{shil}, \textquotedblleft{}flat stone for grinding on\textquotedblright{}, from Sanskrit \emph{śilā}, \textquotedblleft{}rock\textquotedblright{} (CDIAL 12459), + \emph{noṛā}, \textquotedblleft{}rolling-pin\textquotedblright{}, which CDIAL lists under Sanskrit \emph{lōṭha}, \textquotedblleft{}rolling\textquotedblright{} (CDIAL 11134).
\end{cousinweb}

\begin{grammarlens}[title={What you've built}]
One more word for this chapter.
\end{grammarlens}

\subsection*{Your turn}
\begin{itemize}
\raggedright
  \item \emph{Say it:} \emph{shilnoṛā}
  \item \emph{Say it:} \emph{shilnoṛā}, once more
  \item \emph{Say it:} say \emph{shilnoṛā}
  \item \emph{Recall:} say the Bengali for a duster, then the Bengali for a bonti, then say \emph{shilnoṛā} again
\end{itemize}

\begin{tcolorbox}[breakable,colback=teal!4,colframe=teal!35!black,title={Before you move on}]
What does \bn{শিলনোড়া} mean? (\textquotedblleft{}A grinding stone\textquotedblright{}.) Say it once more.
\end{tcolorbox}
\section[\texorpdfstring{\bn{গ্যাস} (gæs)}{গ্যাস (gæs)}]{\bn{গ্যাস} (gæs) — gas}

\label{lesson:BN-C175-gs}



\begin{quote}
Before the new one: say the Bengali for a bonti, then the Bengali for a grinding stone.
\end{quote}

\subsection*{You'll want to know: \bn{গ্যাস}}
\textbf{\bn{গ্যাস}} — \emph{gæs} — \textquotedblleft{}gas\textquotedblright{}.

\begin{cousinweb}
\textbf{\bn{গ্যাস}} (\emph{gæs}) is a loan from English \emph{gas}.
\end{cousinweb}

\begin{grammarlens}[title={What you've built}]
One more word for this chapter.
\end{grammarlens}

\subsection*{Your turn}
\begin{itemize}
\raggedright
  \item \emph{Say it:} \emph{gæs}
  \item \emph{Say it:} \emph{gæs}, once more
  \item \emph{Say it:} say \emph{gæs}
  \item \emph{Recall:} say the Bengali for a bonti, then the Bengali for a grinding stone, then say \emph{gæs} again
\end{itemize}

\begin{tcolorbox}[breakable,colback=teal!4,colframe=teal!35!black,title={Before you move on}]
What does \bn{গ্যাস} mean? (\textquotedblleft{}Gas\textquotedblright{}.) Say it once more.
\end{tcolorbox}
